\section{Introducción}
Esta memoria documenta el desarrollo de la Práctica 1, cuyo objetivo cardinal reside en el diseño e implementación de modelos de aprendizaje automático capaces de predecir el comportamiento financiero a nivel particular. Específicamente, nos enfocamos en predecir el volumen de gastos futuros basándonos en variables históricas y comportamientos financieros estacionales.

Para el desarrollo del presente proyecto es imperativo aclarar que los datos que se han analizado son estrictamente sintéticos. Durante la simulación de este conjunto de datos hemos intentado que los registros pertenezcan y representen de manera fidedigna el comportamiento económico de una \textbf{sola persona}. La decisión de aislar la información a un nivel individual se fundamenta en un requisito de diseño para el futuro del producto: el modelo final en producción se irá reentrenando y ajustando dinámicamente con los datos exclusivos de cada usuario, garantizando así una alta personalización predictiva y sin mezclar los entrenamientos con los datos financieros de diferentes individuos, preservando tanto la integridad del perfil de gasto como la privacidad intrínseca por usuario.
